\section{Results}

\begin{frame}{Value iterations}

\begin{itemize}
    \item Convergence to $\pi_*$ after 17 iterations ($\sim 8$h on a laptop)
    \item Optimal policy $\pi_*$ plays perfectly against an optimal opponent
\end{itemize}

\end{frame}

\begin{frame}{Why no convergence?}
  \begin{itemize}
    \item Difficulties:
    \begin{itemize}
      \item too many possible actions at each step
      \item ideal opponent that never makes mistakes
      \item very sparse positive rewards in training episodes
    \end{itemize}
    \item Possible solutions:
    \begin{itemize}
      \item more training time 
      \item implement draw by repetition and 50-move rule $\rightarrow$ more frequent terminal states
    \end{itemize}
  \end{itemize}
\end{frame}

\begin{frame}{Policy interpretability through human chess principles}
  
\end{frame}

\begin{frame}{Video Animation of optimal ply vs our policy}
  maybe for a mate in 5 or more

\end{frame}

\begin{frame}
  Future Works

  * zobrist hashing and invariance properties in chess endgames to reduce state space cardinality
\end{frame}

\begin{frame}
Why is it so difficult to solve chess endgames?

\begin{quote}
A player can sometimes afford the luxury of an inaccurate move, or even a definite error, in the opening or middlegame without necessarily obtaining a lost position. In the endgame … an error can be decisive, and we are rarely presented with a second chance.
\end{quote}

\hfill - Paul Keres
\end{frame}